\chapter{Data}
Before we want to know how to visualize data, we may want to know about certain characteristics of data(sets) to exploit in a resulting diagram.

\section{Dataset and Data Types}
Generally, datasets are in the shape of three main types.

\begin{description}
    \item[Tables] Tables are (ordered) collections of rows which contain cells with individual values.
        Each dimension of a row is called an \emph{attribute}.
    \item[Networks] Networks or graphs, which includes trees, are collections of (valued) nodes or \emph{vertices} and (valued) edges between nodes.
        Edges can be directed and undirected.
        A tree is a strict directed acyclic graph.
    \item[Spacial data] Spacial data comes in the form of fields, which are continuous, and geometry, which is spacial.
        An example of continuous spacial data are meshes generated by an MRI scan or a vector field reporting wind speeds.
        Geometric spacial data are explicit spacial positions, such as state borders.
\end{description}

\noindent
The values \emph{themselves} in a dataset can have different kinds of attributes.

\begin{description}
    \item[Categorical] Categorical data can not be measured or counted in the form of numbers.
        Examples are texts, shapes, or items.
        This data type can be further split into two more distinct types.

        \begin{description}
            \item[Nominal] Nominal data does not have an implicit order.
                Examples are eye colour, nationality, or gender.
            \item[Ordinal] Ordinal data has some underlying unarithmetic order.
                Examples are letter grades, education level, or sheet notes.
        \end{description}

    \item[Quantitative] Quantitative data can be expressed in numerical values, making it countable and including statistical data analysis.
        These kinds of data are also known as numerical data.

        Again, this data type can be further split into two more distinct types.

        \begin{description}
            \item[Discrete] Discrete data contains the values that fall under integers or whole numbers.
                The total number of students in a class is an example of discrete data.
                These data can’t be broken into decimal or fraction values.

            \item[Continuous] In contrary, continuous data is in the form of fractional numbers.
                It can e.g., be the height of a person, the length of an object, or income.
        \end{description}
\end{description}

\noindent
An important understanding behind each attribute is the semantics associated with it.
To visualize data, one must keep the meaning behind attributes in mind.
This also includes the occasional step to derive new attributes from given attributes, e.g., calculating the trade balance by subtracting the imports from exports.
